% chuong4.tex — sửa nội dung chương ở đây
\section{Thách thức khi áp dụng pháp luật}\label{sec:4}

Chương 3 cho thấy các hành vi của taphoammo.net đã bị pháp luật cấm từ trước khi sàn hoạt động. Chương này trả lời câu hỏi nghiên cứu thứ hai: vì sao luật đã có mà sàn vẫn tồn tại khoảng bốn năm. Nhóm chia các thách thức thành ba nhóm: kỹ thuật và bảo mật, pháp lý và quản trị. Đây là phân tích của nhóm dựa trên dữ kiện đã công bố, không phải kết luận rằng một cơ quan cụ thể đã xử lý chưa đúng.

\subsection{Thách thức kỹ thuật và bảo mật}\label{sec:4.1}

Quy mô lớn và dữ liệu phân tán khiến việc phát hiện và xác minh khó thực hiện hoàn toàn bằng tay.

\begin{itemize}
\item \textbf{Khối lượng lớn.} Sàn có hơn 46.000 gian hàng và trên 53 triệu lượt giao dịch~\cite{bca-2026}. Đây là số liệu quy mô, không phải số tin vi phạm, nhưng cho thấy việc đọc và kiểm tra từng tin bằng tay là không khả thi.
\item \textbf{Cách viết lách bộ lọc.} Từ lóng, viết tắt và các biến thể như “d.a.t.a” hay “0tp” khiến bộ lọc dựa trên từ khóa đơn lẻ dễ bỏ sót hoặc báo nhầm. Đây là tình huống nhóm giả định khi xây dựng demo ở chương 5, chưa phải dữ kiện đã được chứng minh trong vụ việc.
\item \textbf{Tài khoản đổi chủ.} Một tài khoản đã xác thực có thể được bán và chuyển sang người khác dùng, nên việc xác thực khi tạo tài khoản chưa đủ bảo đảm danh tính người dùng về sau.
\item \textbf{Dữ liệu dễ sao chép.} Dữ liệu cá nhân có thể được sao chép và chuyển qua nhiều bên, nên rất khó truy ra nguồn rò rỉ và thu hồi mọi bản sao.
\item \textbf{Hạ tầng ở nước ngoài.} Cơ quan điều tra nêu việc máy chủ đặt ở nước ngoài gây khó khăn cho truy vết. Bản thân việc đặt máy chủ ở nước ngoài không chứng minh vi phạm, nhưng làm tăng nhu cầu phối hợp với nhà cung cấp hạ tầng và bảo toàn chứng cứ.
\end{itemize}
\subsection{Thách thức pháp lý}\label{sec:4.2}

Có dấu hiệu vi phạm chưa đủ để xử lý; cần xác định đúng chủ thể, đúng thời điểm và đủ chứng cứ.

\begin{itemize}
\item \textbf{Nhiều văn bản, hiệu lực khác nhau theo thời gian.} Một hành vi có thể liên quan cùng lúc đến bảo vệ dữ liệu, an ninh mạng, thương mại điện tử và thanh toán. Như mục~\ref{sec:2.2} đã trình bày, khung pháp luật thay đổi hai lần trong thời gian sàn hoạt động, tạo thành ba giai đoạn, nên mỗi giai đoạn phải áp đúng văn bản có hiệu lực. Luật mới không được dùng để xử lý hành vi xảy ra trước đó.
\item \textbf{Ranh giới trách nhiệm giữa người bán và sàn.} Trách nhiệm của sàn phụ thuộc vào vai trò kiểm soát, việc duyệt tin, việc thu phí và cách xử lý phản ánh. Nghĩa vụ kiểm soát thông tin trên sàn đã có từ trước các luật mới~\cite{vbhn-tmdt-2020}. Với hành vi mua bán dữ liệu, mức phạt tối đa 10 lần theo Luật Bảo vệ dữ liệu cá nhân 2025 tính trên khoản thu từ hành vi vi phạm, không phải 10 lần toàn bộ doanh thu doanh nghiệp, và còn phụ thuộc điều kiện áp dụng~\cite{luat-bvdlcn-2025}.
\end{itemize}
Ngoài ra, quyết định khởi tố chưa phải kết luận có tội. Cho đến khi có bản án có hiệu lực, các bị can vẫn được suy đoán vô tội.

\subsection{Thách thức quản trị}\label{sec:4.3}

Phát hiện chỉ có ý nghĩa khi có đầu mối chịu trách nhiệm và hồ sơ đủ để xử lý.

\begin{itemize}
\item \textbf{Vỏ bọc doanh nghiệp.} Sàn đứng tên một công ty có đăng ký, nhưng có đăng ký doanh nghiệp không có nghĩa mọi hàng hóa và giao dịch trên sàn đều hợp pháp. Kiểm tra danh tính người bán và kiểm tra hàng hóa là hai bước riêng.
\item \textbf{Dữ liệu nằm ở nhiều bên.} Thông tin về tin đăng, thanh toán và hạ tầng nằm ở sàn, ngân hàng hoặc ví điện tử, nhà cung cấp máy chủ và các nền tảng quảng cáo. Cần thống nhất đầu mối và cách bàn giao; bảo toàn nguồn, thời điểm, nhật ký và lịch sử chuyển giao; chia sẻ dữ liệu đúng phạm vi thẩm quyền.
\end{itemize}
Từ chín thách thức này, chương 5 đề xuất sáu giải pháp.

